% =============================================================================
% cap05.tex — PRE-DIMENSIONAMENTO
% =============================================================================

\chapter{PRÉ-DIMENSIONAMENTO}\label{cap:predimensionamento}

Este capítulo estabelece as dimensões dos elementos: a geometria da viga de
duplo-bisel, as seções dos pilares e o arranjo e as diagonais dos sistemas de
contraventamento. Apresenta o método, as ferramentas computacionais empregadas, a
varredura paramétrica que conduziu à geometria adotada e a síntese das dimensões
e dos critérios que as governaram.

\section{MÉTODO ADOTADO}\label{sec:metodo-predim}

O pré-dimensionamento de uma estrutura como a deste trabalho não se resolve por
uma sequência fechada de fórmulas. As dimensões da viga condicionam o peso
próprio, que realimenta os esforços; a largura da seção condiciona a largura do
pilar e o espaço disponível para a ligação de base; e o espaçamento das terças
condiciona o comprimento efetivo de flambagem lateral da viga. O problema é
circular, e a via praticável é a varredura: percorre-se um conjunto de
geometrias candidatas, submete-se cada uma à cadeia completa de verificações e
observa-se o que governa em cada região da grade.

O procedimento adotado difere, num ponto essencial, do que é usual em
pré-dimensionamento. Em vez de avaliar as candidatas por um modelo simplificado,
mais rápido porém distinto daquele que produz o memorial, construiu-se um
mecanismo que executa, para cada ponto da grade, exatamente o mesmo
\textit{pipeline} de verificação empregado no dimensionamento final. Cada geometria é avaliada pelas
\nVerifExt{} verificações descritas no \autoref{cap:dimensionamento}, com o mesmo
modelo estrutural, o mesmo gerador de combinações e os mesmos integradores.

\subsection{Ferramentas computacionais}\label{ssec:ferramentas}

O programa de cálculo, descrito no \autoref{ap:codigo-fonte}, foi escrito em
Python \cite{python:3.12}. A \autoref{tab:software} relaciona as bibliotecas e
ferramentas empregadas, com as versões do ambiente em que os resultados desta
monografia foram produzidos.

\tabela{Software empregado}{\footnotesize
\begin{tabular}{p{4.2cm}lp{6.4cm}}
\toprule
    \textbf{Ferramenta} & \textbf{Versão} & \textbf{Uso} \\
\midrule
    Python \cite{python:3.12}            & 3.12.7 & Linguagem do programa de cálculo \\
    NumPy \cite{numpy:2020}              & 2.4.6  & Cálculo vetorial: geometria, diagramas e integração \\
    SciPy \cite{scipy:2020}              & 1.17.1 & Integração numérica e solução de equações \\
    pandas \cite{pandas:2010}            & 3.0.3  & Tabelas de resultados e exportação do memorial \\
    openpyxl \cite{openpyxl}             & 3.1.5  & Gravação da planilha de resultados \\
    Matplotlib \cite{matplotlib:2007}    & 3.11.0 & Gráficos, diagramas de esforços e desenhos cotados \\
    OpenSeesPy \cite{openseespy:2018}    & 3.8.0.0 & Modelo de elementos finitos de validação \\
    PyVista \cite{pyvista:2019} e VTK \cite{vtk:2006} & 0.48.4 e 9.6.2 & Visualização tridimensional do modelo \\
    pytest \cite{pytest}                 & 9.1.0  & Suíte de testes automatizados \\
    TeX Live \cite{texlive:2023} e abnTeX2 \cite{abntex2:2018} & 2023 e 1.9.7 & Composição da monografia \\
\bottomrule
\end{tabular}}{Elaborado pelo autor (2026)}{software}{}{}

Declara-se o uso de ferramenta de inteligência artificial. Os modelos de linguagem
Claude Opus 5 e Claude Opus 5.5, da Anthropic \cite{claude:2026}, acessados por
meio do Claude Code em 2026, foram empregados na complementação e no ajuste do
código do programa de cálculo, escrito inicialmente pelo autor. Toda a concepção
do trabalho, isto é, o arranjo estrutural, os modelos de cálculo, os critérios de
verificação e as decisões de projeto, é do autor, que revisou as contribuições da
ferramenta e responde pelos resultados.

%% TODO: verificar — conferir a regra da UFPR vigente sobre declaração de uso de
%% inteligência artificial em trabalhos acadêmicos (local e forma da declaração).

\section{VARREDURA GLOBAL DA VIGA}\label{sec:varredura-global}

\subsection{Grade percorrida}\label{ssec:grade}

Percorreram-se as três variáveis que definem a conformação da viga de
duplo-bisel: a inclinação do extradorso, entre 6,0$^{\circ}$ e 16,0$^{\circ}$ em
incrementos de 0,5$^{\circ}$; a altura no apoio, entre 0,40 m e 1,20 m em
incrementos de 0,05 m; e a largura da seção, entre 0,20 m e 0,36 m em
incrementos de 0,02 m. A largura dos pilares acompanha a da viga em cada ponto,
por decisão de projeto quanto ao alinhamento do apoio, de modo que a varredura
não avalia estruturas que o projeto não construiria. A \autoref{fig:esq-variaveis}
identifica as variáveis na viga: a inclinação $\alpha$ e a altura no apoio $h_0$
definem, para o vão $L$, a altura na cumeeira $h_{ap}$, e a largura $b$ é constante
ao longo do vão.

\figura
{Variáveis da varredura global da viga de duplo-bisel}
{1.0}
{fig/esq01_variaveis_viga.pdf}
{Elaborado pelo autor (2026)}
{esq-variaveis}
{}
{}

Cada geometria inclui o balanço do beiral, de 1,135 m além do apoio externo, e
é avaliada pela cadeia completa de verificações, inclusive as ligações, as bases,
a calha central e os pilares, cuja largura acompanha a da viga. Resultaram 3.213
geometrias, das quais 1.111 foram aprovadas em todas as verificações e 2.102
reprovadas. As reprovações vêm de três regiões da grade, e só na última intervém a
resistência da viga.

\begin{alineas}
    \item Todas as geometrias com largura de até 0,28 m são reprovadas. Em todas,
          reprova a ligação entre terça e viga, cujos parafusos inclinados exigem
          largura da viga para a penetração e as distâncias às bordas. Com 0,20 e
          0,22 m reprovam também as bases dos pilares e o incêndio de veículo e,
          com 0,20 m, o pilar central, cuja taxa média sobe de 0,40 com
          $b = 0{,}32$ m para 0,63 com $b = 0{,}24$ m e 1,05 com $b = 0{,}20$ m.
    \item Todas as geometrias com inclinação abaixo de 8,0$^{\circ}$ são reprovadas
          pela calha central da \autoref{ssec:calha-central}. Com a borda da
          calha 232 mm acima da terça T9, o último pano de telha fica com
          inclinação menor que os 5~\% exigidos pela telha. Nas de menor altura
          no apoio, reprova também o abaixamento da cumeeira, que passa de 5~\%
          da flecha do eixo (\autoref{sec:ordem-analise}).
    \item Com largura de 0,30 m ou mais, reprovam apenas as geometrias de menor
          altura no apoio: com 0,40 m, pelo cisalhamento e pelas tensões na
          borda inclinada (com 8,0$^{\circ}$ e 0,30 m, também pelo abaixamento da
          cumeeira), e, com inclinação acima de 13$^{\circ}$, também com
          0,45 m, pela borda inclinada. Além dessas, reprovam só quatro das geometrias mais pesadas da grade,
          com largura de 0,36 m, inclinação de 15$^{\circ}$ ou mais e altura no
          apoio de 1,15 m ou mais, e reprova o sistema, não a viga:
          a lâmina do pilar central montante soma ao esforço do painel vertical a
          reação de duas vigas mais pesadas. O $\alpha_{cr}$ transversal, calculado
          com a rigidez do diafragma da cobertura (\autoref{sec:ordem-analise}),
          fica acima de 10 em todas as geometrias aprovadas.
\end{alineas}

\subsection{Resultado}\label{ssec:resultado-varredura}

A \autoref{fig:varredura-global} apresenta o espaço percorrido, com o volume de
madeira lamelada colada das vigas principais no eixo das abscissas e, no eixo
das ordenadas, a maior taxa de utilização entre as verificações que dependem da
seção da viga. Cada ponto corresponde a uma execução completa do
\textit{pipeline} de verificação; a taxa representada é a maior entre
cisalhamento, flexo-compressão, cumeeira, borda inclinada, flambagem lateral,
apoio e estado-limite de serviço, e o volume refere-se às 18 vigas
principais da edificação.

\figura
{Espaço de geometrias percorrido na varredura}
{0.98}
{fig/varredura_global.pdf}
{Elaborado pelo autor (2026)}
{varredura-global}
{}
{}

A figura torna visível o resultado central desta etapa. Na fronteira inferior da
região aprovada, a taxa da viga cai rapidamente com o volume. Acima de cerca de
170 m\textsuperscript{3}, ela entra num patamar entre 0,5 e 0,6, no qual o
acréscimo de madeira já não a reduz. A geometria adotada situa-se nesse patamar.

A \autoref{tab:fronteira-volume} quantifica esse patamar. Cada linha é a
geometria de menor volume, entre as 1.111 aprovadas, cuja maior taxa de
utilização na viga não excede a margem indicada. Nenhuma geometria aprovada tem
taxa na viga abaixo de 0,48, e por isso a tabela para na margem de 0,50. As duas
últimas colunas registram as taxas do contraventamento vertical e do pilar de
canto na mesma geometria, obtidas pela EN 1995-1-1 \cite{EN1995-1-1:2004} e pela
EN 1992-4 \cite{EN1992-4:2018}. A tabela e a figura são geradas pelo programa a
partir do resultado da varredura.

% tab_fronteira_volume.tex — GERADA por mlc_projeto/predim/varredura_global.py (issue 93): 1111 aprovadas de 3213. NÃO EDITAR À MÃO.
\tabela{Geometria de menor volume para cada política de margem na viga}{\footnotesize
\begin{tabular}{ccccrccc}
\toprule
    \textbf{Margem} & \textbf{$\alpha$} & \textbf{$h_0$} & \textbf{$b$} &
    \textbf{Volume} & \textbf{Verificação} & \textbf{Contrav.} & \textbf{Pilar} \\
    \textbf{na viga} & \textbf{(graus)} & \textbf{(m)} & \textbf{(m)} &
    \textbf{(m\textsuperscript{3})} & \textbf{governante} & \textbf{vertical} & \textbf{de canto} \\
\midrule
    $\eta \leq 1,00$ & 8,0 & 0,45 & 0,30 & 129,3 & Cisalhamento & 0,79 & 0,81 \\
    $\eta \leq 0,90$ & 8,0 & 0,50 & 0,30 & 135,1 & Cisalhamento & 0,79 & 0,81 \\
    $\eta \leq 0,75$ & 8,0 & 0,60 & 0,30 & 146,6 & Cisalhamento & 0,80 & 0,82 \\
    $\eta \leq 0,60$ & 8,0 & 0,70 & 0,30 & 158,1 & Cisalhamento & 0,81 & 0,83 \\
    $\eta \leq 0,50$ & 8,0 & 0,70 & 0,36 & 189,8 & Apoio central & 0,82 & 0,84 \\
\midrule
    \textbf{Adotada} & \textbf{10,0} & \textbf{1,00} & \textbf{0,32} & \textbf{226,6} & \textbf{Apoio central} & \textbf{0,84} & \textbf{0,89} \\
\bottomrule
\end{tabular}}{Elaborado pelo autor (2026)}{fronteira-volume}{}{}


\subsection{O que governa o dimensionamento}\label{ssec:o-que-governa}

A leitura conjunta da figura e da tabela conduz a três constatações.

A primeira é que, na geometria adotada, as maiores taxas calculadas não estão
na viga. A maior entre as verificações da viga é a do apoio sobre a sela
central, com 0,58, seguida do apoio sobre o pilar lateral, com 0,49, do
cisalhamento, com 0,41, e das tensões na zona de cumeeira, com 0,36. As maiores
taxas estão nos componentes de aço e nas ligações, e nos contraventamentos
vertical e da cobertura, com 0,84 e 0,80. Esse contraste resulta em parte das
políticas adotadas: a viga foi escolhida com volume bem acima do mínimo aprovado
(\autoref{sec:geometria-adotada}), e os componentes de aço são dimensionados até
a taxa-alvo de 0,90, como as bases ancoradas e as soldas. Na fronteira de menor
volume da \autoref{tab:fronteira-volume}, ao contrário, é a viga que governa,
pelo cisalhamento. As solicitações dos componentes de estabilização respondem à
área exposta ao vento e às reações de apoio, e são praticamente insensíveis à
altura no apoio e à largura da seção.

A segunda é que os limites da região aprovada são impostos, em parte, por
elementos que acompanham a viga, nos detalhes mantidos fixos na varredura: a
ligação entre terça e viga, a calha central e os perfis comerciais. Outros
detalhes deslocariam a fronteira. A largura mínima de 0,30 m vem da ligação
entre terça e viga. A inclinação
mínima de 8,0$^{\circ}$ vem da calha central. Por isso, na \autoref{tab:fronteira-volume},
as geometrias de menor volume têm todas $\alpha = 8{,}0^{\circ}$ e, até a margem
de 0,60, $b = 0{,}30$ m, e diferem só pela altura no apoio; para descer à margem de
0,50, a largura sobe a 0,36 m. A taxa dos pilares montantes bipartidos
(\autoref{ssec:pilar-bipartido}) varia pouco, entre 0,85 e 1,00, porque as
lâminas têm seção própria e a maior parte do seu esforço vem do painel vertical,
que não depende da viga. A do
contraventamento vertical fica entre 0,78 e 0,89 em todas as geometrias aprovadas,
porque o programa escolhe, em cada uma, o menor diâmetro comercial com taxa até
0,90. O deslocamento transversal do diafragma da cobertura fica entre 0,67 e 0,75
do limite, porque o programa escolhe, em cada geometria, o menor diâmetro das
diagonais dos frontões com taxa de deslocamento até 0,75.

A terceira é que, dentro da região aprovada, a altura no apoio é a variável que
controla o volume e a margem da viga. Com $\alpha = 8{,}0^{\circ}$ e
$b = 0{,}32$ m, o volume cresce de 131,7 m\textsuperscript{3} com $h_0 = 0{,}40$ m
para 205,6 m\textsuperscript{3} com $h_0 = 1{,}00$ m. A taxa da viga cai de 0,98
para 0,53 até $h_0 = 0{,}75$ m, governada pelo cisalhamento e pela borda
inclinada. Daí em diante volta a crescer devagar, até 0,56 com $h_0 = 1{,}00$ m,
porque passa a governar o apoio sobre a sela central, cuja reação cresce com o
peso próprio da viga.

\section{GEOMETRIA ADOTADA}\label{sec:geometria-adotada}

Adotou-se $\alpha = 10{,}0^{\circ}$, $h_0 = 1{,}00$ m e $b = 0{,}32$ m, do que resulta
altura na cumeeira de 2,785 m e volume de 226,6 m\textsuperscript{3} de madeira
lamelada colada nas vigas principais, incluídos os balanços do beiral.

A varredura mostra que existem geometrias aprovadas com cerca de 60~\% desse
volume. A adoção de uma seção mais robusta que o mínimo apoia-se em quatro
razões, que a própria varredura permite sustentar.

A primeira é a largura. Com o tirante interno (\autoref{ssec:topo-pilar}), o pé
dos pilares correntes deixou de levar tração por cavilhas, e a largura de 0,30 m
também aprova em quase toda a grade. Na geometria adotada, porém, ela levaria os
dois modos frágeis da ligação L1, que não admitem redistribuição, para perto do
limite de 0,75: o fendilhamento e os pinos para 0,73, frente a 0,67 e 0,69
com 0,32 m. A economia seria de 14 m\textsuperscript{3}, 6~\% do volume das
vigas. Manteve-se 0,32 m por essa margem nos modos frágeis do apoio da peça
principal. A largura
acompanha a do pilar, o que elimina excentricidade na transferência da reação
vertical, e oferece às ligações a espessura necessária para os grupos de pinos e
de parafusos com as distâncias mínimas exigidas.

A segunda é a altura no apoio. O valor de 1,00 m acomoda a ligação L1, com o rasgo
da chapa embutida, os furos ovalizados e os 100 mm de madeira inteira acima do
rasgo (\autoref{ssec:altura-rasgo}), e dá ao balanço do beiral altura de 0,80 m na
ponta. As geometrias de menor volume operam com alturas no apoio entre 0,40 m e
0,65 m, faixa em que a ligação passa a competir por espaço com a própria seção e
em que a taxa da viga se aproxima de 1,0.

A terceira é a margem na viga. Com $h_0 = 1{,}00$ m, a maior taxa da viga é 0,58,
de modo que a viga, peça principal e de substituição inviável, não está entre os
elementos que governam a segurança do conjunto.

A quarta é a inclinação. A inclinação de 10,0$^{\circ}$ fica 2,0$^{\circ}$ acima do
mínimo de 8,0$^{\circ}$ imposto pela calha, o que dá ao último pano de telha
inclinação de 9,2~\% em vez dos 5~\% mínimos e mantém a altura na cumeeira em
2,785 m, compatível com o transporte rodoviário da peça inteira. Com
$\alpha = 8{,}0^{\circ}$ e as mesmas $h_0$ e $b$, o volume seria de 205,6
m\textsuperscript{3}, 9~\% menor.

\section{PILARES}\label{sec:predim-pilares}

Os pilares laterais e o central receberam seção única de
0,32 $\times$ 0,27 m. A largura acompanha a da viga, pelo motivo já exposto de
alinhamento do apoio. A profundidade no plano do pórtico é governada pelo pilar
central, que recebe a reação de duas vigas e portanto o dobro do esforço normal
dos laterais: sua taxa de utilização é 0,42, frente a 0,21 do lateral.

Os 14 pilares montantes das baías contraventadas, isto é, os de canto, os
laterais P4, P6, P22 e P24 e os centrais montantes, constituem exceção: são
bipartidos, com duas lâminas de 0,16 $\times$ 0,32 m dos lados da viga, que
sobem até a escora. A razão é o nó do painel de contraventamento vertical, que
precisa ficar no nível da escora, e não no topo do pilar, como exposto na
\autoref{ssec:no-topo}; a lâmina leva a componente vertical da diagonal paralela
às fibras. Os de canto acumulam ainda o esforço do painel do frontão.

\section{TERÇAS}\label{sec:predim-tercas}

O espaçamento das terças é fixado em $L/8 = 2{,}53$ m, porque define os oito
painéis por água do contraventamento da cobertura e o comprimento de flambagem
lateral da viga. A seção é única para todas as terças, inclusive as de ponta,
as de beiral, que funcionam como corda do diafragma, e as que trabalham como
escora nas baías contraventadas, para não multiplicar os tipos de peça na
fabricação. Ela é obtida pelo mesmo procedimento da viga: para cada seção
comercial candidata, com larguras de 115 a 200 mm e alturas múltiplas de 40 mm,
de 240 a 600 mm, a cadeia completa de verificações é executada, já que a seção
altera o peso próprio, e portanto todas as ações, e as ligações que a
atravessam. Aceita-se a candidata que resulta aprovada e em que cada verificação
dependente da terça tem taxa não superior a 0,75, a mesma folga de 25\% adotada
para os pilares; adota-se a de menor área.

\tabela{Seções de terça candidatas em torno da adotada}{\footnotesize
\begin{tabular}{lrcl}
\toprule
    \textbf{Seção} (m) & \textbf{Área} (m\textsuperscript{2}) & \textbf{$\eta$ máximo} & \textbf{Situação} \\
\midrule
    0,16 $\times$ 0,44 & 0,0704 & 0,92 & Acima da folga adotada \\
    0,20 $\times$ 0,36 & 0,0720 & 0,83 & Não atende \\
    0,18 $\times$ 0,40 & 0,0720 & 0,76 & Acima da folga adotada \\
    0,18 $\times$ 0,44 & 0,0792 & 0,74 & Atende, menor área \\
    \textbf{0,20 $\times$ 0,40} & \textbf{0,0800} & \textbf{0,74} & \textbf{Adotada} \\
    0,20 $\times$ 0,56 & 0,1120 & 0,73 & Atende, com 40\% a mais de madeira \\
\bottomrule
\end{tabular}}{Elaborado pelo autor (2026)}{varredura-tercas}{}{}

A escora central das baías contraventadas recebe 351 kN e trabalha à
flexocompressão, com a flexão da cobertura. Com a terça corrente, de 200 mm de
largura, a flambagem em torno do eixo de menor inércia levaria a sua taxa a 1,03,
e por isso ela tem seção própria, de 0,24 por 0,40 m, apenas nas quatro baías
contraventadas do alinhamento central, com taxa de 0,65. Nas demais candidatas
governam a ligação das escoras ao nó e a ligação entre terça e viga, todas
próximas de 0,74. A seção de 0,18 por 0,44 m tem área 1\% menor que a adotada e
também atende; manteve-se a de 0,20 por 0,40 m, pela diferença desprezível e por
conservar a largura das escoras de beiral e das juntas já detalhadas. A terça
como barra fica longe do limite: a interação de flexão e compressão da biela do
diafragma tem taxa de 0,59 e a flexão, de 0,20.

\section{DIAGONAIS DE CONTRAVENTAMENTO}\label{sec:predim-contraventamento}

\subsection{Contraventamento da cobertura}\label{ssec:predim-cv-cobertura}

As diagonais do contraventamento horizontal da cobertura cumprem duas funções
simultâneas: resistem à força de estabilização prescrita no item 9.2.5 da
EN 1995-1-1 \cite{EN1995-1-1:2004} e validam as terças como pontos de contenção
lateral da viga. Essa segunda função é a mais restritiva, porque dela depende o
comprimento efetivo de flambagem lateral: se o sistema de cobertura não for
capaz de estabilizar, as terças deixam de ser contenções laterais, e o
comprimento efetivo passa do espaçamento entre terças, de 2,53 m, ao vão da
própria viga, de 20,25 m.

Cada baía contraventada funciona como uma treliça no plano da cobertura,
paralela às vigas: os banzos são as duas vigas que delimitam a baía, os
montantes são as terças e as diagonais são as cruzes. A treliça vence o vão da
viga, entre o beiral e a linha central, e é esse o vão $\ell$ do item 9.2.5.3
da EN 1995-1-1 \cite{EN1995-1-1:2004}, igual a 20,25 m, ao qual corresponde
$k_{\ell} = \min(1;\ \sqrt{15/20{,}25}) = 0{,}86$. Com duas baías rígidas, as
nove vigas de cada água repartem-se entre dois sistemas, cada um estabilizando
$n = 5$ vigas. Para que a treliça exista, as cruzes ocupam os oito painéis do
vão, e não apenas os painéis junto aos apoios; a diagonal mais solicitada é a
do painel de extremidade, onde o cortante da treliça é máximo.

A varredura do diâmetro da diagonal, mantido o arranjo de duas baías rígidas nas
extremidades, consta da \autoref{tab:varredura-diagonal-cobertura}. A
verificação segue o item 9.2.5 da EN 1995-1-1 \cite{EN1995-1-1:2004}, as barras
são redondas em aço S235, e a coluna de estabilização indica se as terças
permanecem válidas como contenção lateral da viga.

\tabela{Varredura do diâmetro da diagonal do contraventamento da cobertura}{\footnotesize
\begin{tabular}{cccl}
\toprule
    \textbf{Diâmetro} & \textbf{$\eta$ da diagonal} & \textbf{Estabiliza} & \textbf{Situação} \\
    \textbf{(mm)} & & \textbf{a cobertura} & \\
\midrule
    28 & 1,32 & Não  & Reprova \\
    30 & 1,15 & Não  & Reprova \\
    33 & 0,95 & Sim  & Acima do alvo de 0,90 \\
    36 & 0,80 & Sim  & \textbf{Adotado} \\
    39 & 0,68 & Sim  & Atende, mais pesado \\
\bottomrule
\end{tabular}}{Elaborado pelo autor (2026)}{varredura-diagonal-cobertura}{}{}

Adotou-se 36 mm, o menor diâmetro com taxa abaixo do alvo de 0,90 das peças de
aço; o programa o escolhe automaticamente, como nas diagonais verticais.

A força na diagonal não é governada pela estabilização das vigas, e sim pelo
vento paralelo à cumeeira, que a cobertura coleta por atrito e conduz às baías
rígidas: a reação de cada apoio da treliça soma 173,8 kN do vento, já majorado
por $\gamma_Q = 1{,}5$, e apenas 1,3 kN da força de estabilização do item 9.2.5.3, que atua em adição às ações
externas. Acrescentar baías rígidas
reduziria ambas as parcelas, porque o vento e as $n$ vigas se repartiriam por
mais sistemas, mas não pouparia aço: a seção de cada diagonal diminuiria na
mesma proporção em que o número de diagonais cresceria, e o número de ligações
aumentaria. Manteve-se, por isso, o arranjo de apenas duas baías rígidas, situadas
nas extremidades, CV-H1 entre os eixos 1 e 2 e CV-H2 entre os eixos 8 e 9, como
mostra a \autoref{fig:des-planta-cv}. As baías intermediárias recebem cruzes leves,
as CV-D, que não servem ao vento longitudinal, mas formam a alma do diafragma pelo
qual a força transversal chega aos frontões; a sua necessidade e o seu
dimensionamento constam da \autoref{sec:ordem-analise}.

\figura
{Planta de contraventamento}
{1.0}
{fig/des02_planta_contraventamento.pdf}
{Elaborado pelo autor (2026)}
{des-planta-cv}
{}
{}

\subsection{Contraventamento vertical}\label{ssec:predim-cv-vertical}

Os painéis de contraventamento vertical conduzem a força longitudinal de vento
do nível da cobertura até a fundação. O número de painéis é a variável de
projeto, e sua escolha não é um problema de resistência da diagonal apenas: cada
painel adicional reduz o esforço descendente no montante, e com ele a solicitação
do pilar de canto e da respectiva ancoragem, mas ocupa uma baía que, em um
estacionamento aberto, deixaria de permitir a circulação de veículos.

Antes do número, decide-se onde colocá-los, e essa decisão decorre do
contraventamento da cobertura. A treliça de cada baía rígida apoia-se no beiral e
no alinhamento central, entregando metade da força em cada apoio. Somadas as duas
águas, cada beiral recebe um quarto da força longitudinal, e o alinhamento
central recebe metade. A força de imperfeição reparte-se na mesma proporção,
porque o pilar central carrega o dobro da carga vertical do pilar de beiral. Se
os painéis ficassem apenas nas fachadas, a reação central das treliças não teria
caminho até a fundação, pois os pilares centrais são articulados no topo e na
base. Há, portanto, dois conjuntos de painéis, nas fachadas e no alinhamento
central, e a variável de projeto é o número de painéis em cada um.

O \textit{pipeline} completo foi executado para sete distribuições, com a
diagonal dimensionada automaticamente para taxa de utilização não superior a
0,90, em aço S235. Nas fachadas, os painéis vão aos pares, um em cada
extremidade de cada fachada; no alinhamento central, um par por extremidade. A
\autoref{tab:distribuicao-paineis} reúne o resultado, com a lâmina mais solicitada
do pilar montante bipartido (\autoref{ssec:pilar-bipartido}) verificada pelo item
6.3.2 da EN 1995-1-1 \cite{EN1995-1-1:2004} e a ancoragem pela EN 1992-4
\cite{EN1992-4:2018}.

\tabela{Distribuição dos painéis entre as fachadas e o alinhamento central}{\scriptsize
\setlength{\tabcolsep}{3pt}
\begin{tabular}{cccccccrcl}
\toprule
    \textbf{Fach.} & \textbf{Central} & \textbf{$\varnothing$ (mm)} & \textbf{$V$ por painel} & \textbf{$\eta$ pilar} & \textbf{$\eta$ pilar} & \textbf{$\eta$ anc.} & \textbf{Aço} & \textbf{Baías} & \textbf{Situação} \\
    & & \textbf{fach./centr.} & \textbf{(kN)} & \textbf{de canto} & \textbf{centr. mont.} & \textbf{central} & \textbf{(kg)} & \textbf{ocup.} & \\
\midrule
     4 & 2 & 45 / 64 & 176 / 351 & 0,88 & 1,54 & 1,06 & 1.919 &  6 & Reprova \\
     8 & 2 & 33 / 64 &  88 / 351 & 0,75 & 1,54 & 1,06 & 1.991 & 10 & Reprova \\
     4 & 4 & 45 / 45 & 176 / 176 & 0,88 & 0,91 & 0,66 & 1.908 &  8 & \textbf{Adotado} \\
     8 & 4 & 33 / 45 &  88 / 176 & 0,75 & 0,91 & 0,66 & 1.980 & 12 & Atende \\
    12 & 4 & 27 / 45 &  59 / 176 & 0,75 & 0,91 & 0,66 & 1.984 & 16 & Atende \\
     4 & 8 & 45 / 33 & 176 / 88  & 0,88 & 0,75 & 0,46 & 1.980 & 12 & Atende \\
     8 & 8 & 33 / 33 &  88 / 88  & 0,75 & 0,75 & 0,46 & 2.052 & 16 & Atende \\
\bottomrule
\end{tabular}}{Elaborado pelo autor (2026)}{distribuicao-paineis}{}{}

Com apenas dois painéis no alinhamento central, cada um recebe 351 kN, a diagonal
passa a 64 mm e o pilar central montante, P2, P5, P8, P20, P23 ou P26, cuja lâmina
soma o esforço descendente de 517 kN à metade das reações das duas vigas, reprova,
qualquer que seja o número de painéis nas fachadas. A partir de quatro painéis centrais,
todas as distribuições atendem.

A varredura revela um resultado que inverte a intuição corrente: o peso total de
aço das diagonais verticais varia pouco com o número de painéis, entre 1.908 e
2.052 kg em toda a faixa examinada. Painéis mais numerosos com
barras mais finas custam o mesmo aço que poucos painéis com barras grossas. O que
o aumento do número de painéis compra é margem no caminho de carga, isto é, no
esforço descendente sobre o montante, na ancoragem e no bloco de fundação. O que
custa são baías de circulação, num estacionamento aberto em que cada painel
fecha uma baía.

Adotaram-se, por isso, quatro painéis em cada conjunto, a menor distribuição que
atende. Nas fachadas, há um painel em cada extremidade de cada fachada, sob as
baías rígidas da cobertura. No alinhamento central, há dois painéis contíguos em
cada extremidade. Cada painel recebe 176 kN, com diagonal de 45 mm, e ficam
livres 12 das 16 baías das fachadas e 4 das 8 do alinhamento central. A
\autoref{fig:des-elevacoes} mostra os painéis nas elevações, com os nomes CV-V1 a
CV-V8.

\figura
{Elevações longitudinais dos alinhamentos A e B}
{1.0}
{fig/des03_elevacoes_longitudinais.pdf}
{Elaborado pelo autor (2026)}
{des-elevacoes}
{}
{}

\subsection{Contraventamento dos frontões}\label{ssec:predim-cv-frontao}

O vento perpendicular às cumeeiras e a força de imperfeição chegam aos frontões
pelo diafragma de cobertura. Em cada frontão, um painel de cruz leva essa força à
fundação. O painel foi dimensionado com largura igual ao espaçamento entre
pórticos, 5,75 m, que dá às diagonais inclinação de 50,6$^{\circ}$ e as mantém
curtas. No plano do frontão, porém, os pilares do pórtico distam 20,25 m entre
si, e o painel precisa de um segundo montante a 5,75 m do canto. Há, por isso,
em cada frontão um pilar de frontão, PF1 no eixo 1 e PF2 no eixo 9, com a mesma seção de 0,32 $\times$ 0,27 m
dos pilares laterais, a 5,75 m do canto. Os painéis CV-T1, entre P1 e PF1, e
CV-T2, entre PF2 e P27, ficam em cantos opostos, de modo que o conjunto é
simétrico em relação ao centro da planta.

O pilar de frontão não apoia a viga. A ligação de topo, L1f, transmite apenas a
força horizontal do painel, por furos oblongos verticais, e deixa a viga fletir
livremente; a base, L3f, recebe o arrancamento e o cortante da cruz. Os dois
pilares acrescentam arrasto de vento nas duas direções, e esse acréscimo foi
incorporado às forças de todos os sistemas de contraventamento. A
\autoref{fig:des-transversal} mostra o pórtico do frontão.

\figura
{Elevação transversal do pórtico do frontão}
{1.0}
{fig/des04_elevacao_transversal.pdf}
{Elaborado pelo autor (2026)}
{des-transversal}
{}
{}

\section{SÍNTESE}\label{sec:geometria-resultante}

A \autoref{tab:criterios-predim} sintetiza as dimensões estabelecidas nesta
etapa e o critério que governou cada uma, conforme a justificativa da
\autoref{sec:geometria-adotada}. As dimensões da viga adotada não foram
governadas pela sua resistência, que deixa margem
(\autoref{ssec:o-que-governa}), e sim pelos limites e pelas margens ali
expostos. O espaçamento das terças é fixado pela
geometria do contraventamento, e a sua seção resulta da varredura da
\autoref{sec:predim-tercas}. As verificações que confirmam todas elas constam do
\autoref{cap:dimensionamento}.

\tabela
{Dimensões adotadas e critério governante de cada uma}
{\footnotesize
\begin{tabular}{llp{5.8cm}}
\toprule
    \textbf{Elemento} & \textbf{Dimensão} & \textbf{Critério governante} \\
\midrule
    Inclinação do extradorso & 10,0$^{\circ}$ & Mínimo de 8,0$^{\circ}$ da calha, com margem de drenagem, e altura de transporte \\
    Altura da viga no apoio & 1,00 m & Ligação L1 com o rasgo e a madeira acima dele, margem na viga e altura do beiral \\
    Altura da viga na cumeeira & 2,785 m & Derivada: $h_{ap} = h_0 + (L/2)\tan\alpha$ \\
    Largura da viga & 0,32 m & Margem nos modos frágeis da L1, alinhamento com o pilar e encaixe da chapa \\
    Pilares lateral e central & 0,32 $\times$ 0,27 m & Esforço normal do pilar central \\
    Montantes (14), 2 lâminas & 0,16 $\times$ 0,32 m & Nó do painel vertical no nível da escora \\
    Pilar de frontão & 0,32 $\times$ 0,27 m & Montante do painel do frontão, a 5,75 m do canto \\
    Terças & 0,20 $\times$ 0,40 m & Varredura das terças e função de escora de beiral \\
    Diagonal da cobertura & $\varnothing$ 36 mm & Vento paralelo à cumeeira coletado pela cobertura e levado às baías rígidas \\
    Painéis de contraventamento vertical & 4 + 4 & Fachadas e alinhamento central; caminho de carga das treliças da cobertura \\
    Painéis dos frontões & 1 por frontão & Um painel CV-T em cada frontão, entre o canto e o pilar de frontão \\
\bottomrule
\end{tabular}}
{Elaborado pelo autor (2026)}
{criterios-predim}
{}
{}
